\begin{figure}[ht]
\centering
\begin{tikzpicture}[x=0.85cm, y=0.5cm, font=\footnotesize]

% Semilog: iteration count (x) vs log10 relative residual (y).
% Gauss-Seidel converges linearly with rate ~ (1 - 1/kappa); slow.
% CG converges with rate ~ (sqrt(kappa)-1)/(sqrt(kappa)+1); faster,
% and terminates by step n in exact arithmetic. Curves drawn for a
% symmetric positive definite system with kappa ~ 100.

% Axes
\draw[->, thick] (0,0) -- (13,0) node[right] {iteration $k$};
\draw[->, thick] (0,-12) -- (0,1) node[above, align=center]
  {$\log_{10}\dfrac{\norm{\vect{r}_{k}}}{\norm{\vect{r}_{0}}}$};

\foreach \k in {0,2,4,6,8,10,12} {
  \draw (\k,0) -- (\k,0.25);
  \node[anchor=south, font=\scriptsize] at (\k,0.25) {\k};
}
\foreach \d in {0,-3,-6,-9,-12} {
  \draw (0,\d) -- (-0.2,\d);
  \node[anchor=east, font=\scriptsize] at (-0.2,\d) {$10^{\d}$};
}

% Gauss-Seidel: linear, rate ~ 0.82 per step (slope ~ -0.086 dec/step)
% drawn over the full window, reaching only ~10^{-1.1} by step 12.
\draw[thick, engred] plot coordinates {
  (0,0) (2,-0.18) (4,-0.36) (6,-0.55) (8,-0.74) (10,-0.93) (12,-1.12)
};
\node[engred, font=\scriptsize, anchor=west] at (7.2,-0.55)
  {Gauss-Seidel};

% Conjugate gradients: faster geometric drop, hitting 10^{-11} by step 12
\draw[thick, engblue] plot coordinates {
  (0,0) (1,-0.6) (2,-1.5) (3,-2.6) (4,-3.8) (5,-5.0) (6,-6.1)
  (7,-7.2) (8,-8.2) (9,-9.1) (10,-9.9) (11,-10.6) (12,-11.2)
};
\node[engblue, font=\scriptsize, anchor=west] at (5.4,-4.4)
  {conjugate gradients};

\end{tikzpicture}
\caption[Gauss-Seidel versus conjugate gradients on an SPD system]{Relative
residual norm against iteration count for a symmetric positive definite
linear system of condition number $\kappa \approx 100$, comparing the
stationary Gauss-Seidel iteration with the Krylov conjugate-gradient (CG)
method. Gauss-Seidel contracts the residual by a near-constant factor per
step (linear convergence with asymptotic rate governed by $1 - O(1/\kappa)$),
losing roughly one decimal digit over twelve iterations. CG contracts at
the faster rate $(\sqrt{\kappa}-1)/(\sqrt{\kappa}+1)$, reaching eleven
digits in the same twelve iterations, and would terminate exactly by step
$n$ in exact arithmetic. The gap widens with $\kappa$: the iteration count
to a fixed accuracy scales as $\kappa$ for the stationary method and as
$\sqrt{\kappa}$ for CG. Source: authors' analysis; the convergence bounds
are standard, see \cite[ch.~38]{text:trefethen-bau1997}.}
\label{fig:vol02:ch16:cg-vs-gs}
\end{figure}
